\chapter{So Sánh Thuật Toán}
\label{chap:algorithm-comparison}

\section{So Sánh Lý Thuyết}

Gọi $b$ là hệ số rẽ nhánh của đồ thị và $d$ là độ sâu của đích nông nhất tính từ điểm xuất phát, $m$ là độ sâu tối đa của cây tìm kiếm.

\begin{table}[h!]
\centering
\caption{So sánh lý thuyết của các thuật toán được triển khai}
\begin{tabular}{lccccc}
\toprule
Thuật toán & Độ Phức Tạp Thời Gian & Sử Dụng Bộ Nhớ & Hoàn Chỉnh & Tối Ưu \\
\midrule
Breadth First Search & $O(b^d)$ & $O(b^d)$ & Yes & Only for equal costs \\
Depth First Search & $O(b^m)$ & $O(bm)$ & With a visited set & No \\
Uniform Cost Search & $O(b^{1 + \lfloor C^*/\epsilon \rfloor})$ & $O(b^{1 + \lfloor C^*/\epsilon \rfloor})$ & Yes & Yes \\
A* Search & $O(b^d)$ & $O(b^d)$ & Có & Có \\
Greedy BFS & $O(b^m)$ & $O(b^m)$ & Không & Không \\
Hill Climbing & $O(\infty)$ & $O(1)$ & Không & Không \\
\bottomrule
\end{tabular}
\end{table}

Đây là các giới hạn trong trường hợp xấu nhất. Đối với A*, khi kết hợp với hàm heuristic tốt, thời gian chạy thực tế và bộ nhớ tiêu thụ sẽ giảm đáng kể so với giới hạn lý thuyết này.

\section{Hiệu Suất Thực Tế Trên Tập Dữ Liệu}

Đã chạy các thuật toán trên cùng một tập hợp điểm xuất phát (DL01) và đích (DL15) từ dữ liệu của nhóm với kết quả đo được như sau (chiến thuật balanced):

\begin{table}[h!]
\centering
\caption{Hiệu suất đo được trên tập dữ liệu dự án}
\begin{tabular}{lcccc}
\toprule
Thuật toán & Chi Phí Đường Đi & Số Đỉnh Đã Duyệt & Thời Gian (ms) & Số Cạnh Của Đường Đi \\
\midrule
Breadth First Search & \fillin{value} & \fillin{value} & \fillin{value} & \fillin{value} \\
Depth First Search & \fillin{value} & \fillin{value} & \fillin{value} & \fillin{value} \\
Uniform Cost Search & \fillin{value} & \fillin{value} & \fillin{value} & \fillin{value} \\
A* Search & 9.33 & 12 & 0.46 & 4 \\
Greedy BFS & 9.33 & 5 & 0.04 & 4 \\
Hill Climbing & 9.33 & 5 & 0.03 & 4 \\
\bottomrule
\end{tabular}
\end{table}

\section{Chất Lượng Lộ Trình, Số Đỉnh Đã Duyệt và Thời Gian Xử Lý}

Dựa vào bảng kết quả đo được, có thể thấy với cặp điểm xuất phát và đích cụ thể này, cả ba thuật toán đều tìm được lộ trình với chi phí đường đi (9.33) và số cạnh (4) hoàn toàn giống nhau. Điều này cho thấy trong một số trường hợp địa hình đồ thị thuận lợi, các thuật toán tìm kiếm cục bộ hoặc tham lam vẫn có khả năng đi thẳng đến đích và tìm ra đường đi tốt nhất. Tuy nhiên, sự khác biệt lớn nhất nằm ở không gian tìm kiếm và tốc độ xử lý. A* đã phải duyệt nhiều đỉnh nhất (12 đỉnh) và tốn nhiều thời gian nhất (0.46 ms) do thuật toán này phải đánh giá toàn diện không gian trạng thái xung quanh dựa trên chi phí thực tế nhằm đảm bảo tính tối ưu tuyệt đối. Trong khi đó, Greedy BFS và Hill Climbing hướng thẳng về đích nhờ hàm heuristic, chỉ duyệt đúng 5 đỉnh nằm trên chính lộ trình cuối cùng với thời gian thực thi cực nhỏ (0.04 ms và 0.03 ms). Dưới đây là một số phân tích sâu hơn về đặc điểm lý thuyết giải thích cho kết quả thực tế này:

\begin{itemize}[itemsep=2pt]
  \item \textbf{Thiếu đánh giá của BFS, DFS và USC}
  \item \textbf{A*:} Luôn đảm bảo tìm ra đường đi có tổng chi phí thấp nhất. So với Greedy hay Hill Climbing, số lượng đỉnh phải duyệt có thể nhiều hơn để đảm bảo tính tối ưu, nhưng vẫn hiệu quả nhờ hàm heuristic.
  \item \textbf{Greedy BFS:} Thường có số đỉnh duyệt ít hơn và thời gian chạy nhanh hơn A*, nhưng chi phí đường đi thường cao hơn vì nó không quan tâm đến $g(n)$, dễ dẫn đến các đường đi vòng.
  \item \textbf{Hill Climbing:} Số lượng đỉnh duyệt cực kì ít (bằng đúng độ dài đường đi nếu không bị kẹt), chạy rất nhanh. Tuy nhiên khả năng tìm được đường đến đích trên đồ thị giao thông thực tế khá thấp vì rất dễ đi vào ngõ cụt.
\end{itemize}

\section{Ảnh Hưởng Của Tình Trạng Kẹt Xe Đến Lộ Trình}

Thực hiện chạy từ đỉnh xuất phát DL01 đến đích DL15 hai lần bằng thuật toán A*, một lần ưu tiên thời gian chạy nhanh (kẹt xe thấp) và một lần tránh kẹt xe (kẹt xe cao).

Một lộ trình được ưu tiên khi đường thông thoáng không phải lúc nào cũng là lộ trình tốt nhất khi xét đến kẹt xe. Hàm chi phí được tính toán bằng cách cộng thêm hình phạt kẹt xe (tùy thuộc vào chiến thuật, ví dụ "avoid\_traffic").

\begin{table}[h!]
\centering
\caption{Ảnh hưởng của tình trạng kẹt xe đến lộ trình đã chọn}
\begin{tabular}{lcc}
\toprule
Trạng Thái & Lộ Trình Đã Chọn & Tổng Chi Phí \\
\midrule
Kẹt xe thấp & DL01 $\rightarrow$ DL22 $\rightarrow$ DL21 $\rightarrow$ DL04 $\rightarrow$ DL15 & 12.94 \\
Kẹt xe cao & DL01 $\rightarrow$ DL22 $\rightarrow$ DL21 $\rightarrow$ DL04 $\rightarrow$ DL15 & 2.93 \\
\bottomrule
\end{tabular}
\end{table}

Trong trường hợp này, mặc dù chiến thuật phạt kẹt xe cao được áp dụng, thuật toán vẫn chọn đi qua cùng một lộ trình. Điều này cho thấy trục đường này có sự vượt trội tuyệt đối về cự ly và thời gian so với các tuyến đường vòng, khiến cho ngay cả khi cộng thêm hình phạt kẹt xe, nó vẫn mang lại chi phí tổng thể thấp nhất so với các ngã rẽ khác.
